\section{مبانی و تاریخچۀ پردازش زبان طبیعی}
پردازش زبان طبیعی، حوزه‌ای میان‌رشته‌ای از علوم کامپیوتر، هوش مصنوعی و زبان‌شناسی است که تحلیل، درک و تولید زبان انسانی را با روش‌های محاسباتی بررسی می‌کند. زبان تنها دنباله‌ای از نشانه‌ها نیست؛ معنای واژه و جمله به ساختار، بافت، دانش زمینه و گاه هدف گوینده وابسته است. ازاین‌رو تبدیل متن به یک نمایش عددی، همواره با انتخاب‌هایی دربارۀ اطلاعات حفظ‌شده و اطلاعات ازدست‌رفته همراه است \cite{ref24,ref25}.

در تحلیل سلامت روان، این مسئله اهمیت بیشتری دارد. واژۀ منفرد، نفی، اشاره به خود یا دیگری و ترتیب جمله می‌توانند تفسیر متفاوتی ایجاد کنند. حضور یک واژۀ مرتبط با خودکشی لزوماً نشان‌دهندۀ افکار شخصی نویسنده نیست؛ نقل قول، خبر، بحث حمایتی یا تجربۀ گذشته نیز ممکن است همان واژگان را داشته باشند. بنابراین نتیجۀ طبقه‌بند باید در دامنۀ برچسب‌گذاری داده و محدودیت بازنمایی آن تفسیر شود.

\subsection{وظایف اصلی و پیوند میان آن‌ها}
درک زبان طبیعی به استخراج معنا و ساختار از ورودی و تولید زبان طبیعی به ساخت خروجی زبانی می‌پردازد. ورودی می‌تواند متن یا گفتار باشد؛ تشخیص گفتار و تشخیص نوری نویسه‌ها، مسیرهای تبدیل ورودی به متن‌اند، نه تعریف تمام سطوح پردازش زبان. برخی وظایف مهم عبارت‌اند از واژه‌سازی، تحلیل صرفی، برچسب‌گذاری نقش دستوری، تحلیل نحوی، تشخیص موجودیت نامدار، رفع هم‌مرجعی، تحلیل گفتمان، ترجمه و خلاصه‌سازی.

خلاصه‌سازی، نمایش کوتاه‌تری از متن با حفظ اطلاعات مورد نظر ایجاد می‌کند. رفع هم‌مرجعی مشخص می‌کند کدام عبارت‌ها به یک موجودیت اشاره دارند؛ برای نمونه، مرجع یک ضمیر ممکن است تنها از جمله‌های قبل روشن شود. تحلیل گفتمان روابط میان جمله‌ها، رویدادها و نقش آن‌ها در انسجام متن را بررسی می‌کند. ترجمۀ ماشینی نیز باید علاوه بر واژه‌ها، ساختار و تفاوت کاربرد زبان‌ها را مدیریت کند.

تحلیل صرفی، ساختار واژه و اجزای آن را بررسی می‌کند. برچسب‌گذاری نقش دستوری به هر توکن نقشی مانند اسم یا فعل نسبت می‌دهد. تشخیص موجودیت نامدار، بخش‌هایی مانند نام فرد، سازمان و مکان را شناسایی می‌کند. این وظایف به هم وابسته‌اند؛ برای مثال شناخت موجودیت می‌تواند به رفع هم‌مرجعی کمک کند و تحلیل ساختار جمله می‌تواند بر تفسیر نفی اثر بگذارد. خط‌لولۀ یک کاربرد الزاماً همۀ این مراحل را اجرا نمی‌کند.

در یک معماری پیمانه‌ای، خروجی هر ابزار به ورودی ابزار بعد تبدیل می‌شود. راسپوچر و همکاران نمونه‌ای از استخراج رویداد میان‌زبانی برای انگلیسی، هلندی و ایتالیایی را بررسی کرده‌اند که پردازش پایه، پیوند موجودیت، نقش‌های معنایی و نرمال‌سازی زمانی را به نمایش دانش رویدادمحور متصل می‌کند \cite{ref26}. چنین معماری‌ای امکان جایگزینی ماژول و پیکربندی خط‌لوله را فراهم می‌کند، اما سازگاری قالب و انباشت خطای مراحل نیز باید بررسی شود.

ابهام می‌تواند واژگانی، نحوی یا وابسته به بافت باشد. یک واژه ممکن است چند معنا داشته باشد و یک جمله بیش از یک ساختار نحوی بپذیرد. روش‌ها می‌توانند ابهام را زود کاهش دهند، چند تفسیر را نگه دارند، با کاربر تعامل کنند یا به تفسیرها وزن احتمالی بدهند \cite{ref27}. هیچ انتخابی برای همۀ کاربردها بهترین نیست؛ نگه داشتن تفسیرهای بیشتر هزینه را افزایش می‌دهد و حذف زودهنگام ابهام می‌تواند اطلاعات لازم را از بین ببرد.

\subsection{سیر تاریخی و تحول روش‌ها}
پژوهش‌های ترجمۀ ماشینی از دهه‌های میانی قرن بیستم، به‌ویژه برای زوج‌زبان‌هایی مانند روسی و انگلیسی، یکی از محرک‌های اولیۀ این حوزه بودند. گزارش \en{ALPAC} در ۱۹۶۶ به کاهش حمایت از بخشی از فعالیت‌های ترجمۀ ماشینی انجامید؛ با این حال، به معنای پایان تمام پژوهش‌های زبان طبیعی نبود. کاربردهای محدودتر و مطالعات زبان‌شناختی و ادبیِ رایانه‌ای ادامه یافتند \cite{ref28,ref29}. سامانه‌های پرسش‌پاسخ و رابط زبان طبیعی برای پایگاه داده، نمونه‌هایی از پژوهش‌های دامنه‌محور این دوره‌اند \cite{ref30,ref31,ref32}.

در دهۀ ۱۹۸۰، دستور زبان محاسباتی، معناشناسی، مدل‌سازی دانش و هدف کاربر گسترش یافتند. سامانه‌هایی مانند \en{Core Language Engine} و نظریۀ بازنمایی گفتمان، ابزارهایی برای تحلیل فراتر از واژه و جملۀ منفرد فراهم کردند \cite{ref33,ref34}. ابزارهای زبانی و ارزیابی‌های برنامه‌های پژوهشی، از جمله \en{Alvey} و \en{ARPA}، به توسعۀ روش‌های قابل مقایسه کمک کردند \cite{ref35,ref36,ref37}. مدل‌سازی کاربر و تولید متن منسجم نیز به شناخت نقش بافت و ساختار گفتمان افزودند \cite{ref38,ref39,ref40}.

در دهۀ ۱۹۹۰، روش‌های آماری، استفاده از پیکره و ارزیابی تجربی جایگاه پررنگ‌تری یافتند. رفع ابهام معنایی، استخراج اطلاعات و خلاصه‌سازی از حوزه‌های متأثر از این تحول بودند \cite{ref41,ref42,ref43}. در ادامه، بازنمایی‌های توزیعی واژه و سپس مدل‌های عصبی زمینه‌مند امکان یادگیری ویژگی را افزایش دادند \cite{ref78,ref81,ref82}. این تحول، نیاز به جداسازی درست دادۀ آموزش و ارزیابی را کاهش نداد؛ هر بخشی که از پیکره چیزی یاد می‌گیرد می‌تواند منشأ نشت اطلاعات باشد.

\subsection{از چارچوب عمومی تا آزمایش این پژوهش}
در این پایان‌نامه، واژه‌سازی ساده و میانگین بردار \en{Word2Vec} انتخاب شده است. این انتخاب تحلیل نحوی یا درک کامل معنا را اجرا نمی‌کند؛ هدف آن ساخت یک بازنمایی فشرده برای مقایسۀ طبقه‌بندها و آموزش مستقل مدل‌های محلی است. روش‌های \en{TF-IDF}، \en{Sentence-BERT} و مدل‌های زمینه‌مند، در جایگاه پیشینه و گزینه‌های جایگزین مطرح‌اند و معیارهای آزمایش حاضر به آن‌ها نسبت داده نمی‌شوند.

شکل~\ref{fig:nlp-workflow-corrected} ترتیب اجرایی نسخۀ اصلاح‌شده را نشان می‌دهد. جایگاه تقسیم داده و یادگیری بازنمایی عمداً روشن است: پیکرۀ آزمون به آموزش واژگان و انتخاب مدل وارد نمی‌شود. ابزار دنا نیز از این مسیر اصلی جداست. بخش‌های بعد، تعریف پاکسازی، ابعاد بازنمایی، فضای جست‌وجوی مدل‌ها و قاعدۀ انتخاب آستانه را به‌صورت دقیق ارائه می‌کنند.

\begin{figure}[p]
\centering
\begin{tikzpicture}[
  node distance=7mm,
  stage/.style={draw=blue!45!black,fill=blue!3,rounded corners=3pt,
    inner sep=7pt,font=\small},
  flow/.style={-{Latex[length=2.3mm]},thick,draw=blue!50!black}]
\node[stage] (raw) {\diagramtext{مجموعۀ خام ۲۳۲٬۰۷۴ ردیفی\\کنترل برچسب و متن، حذف تکرار خام و نگهداری نسخۀ داده}};
\node[stage,below=of raw] (split) {\diagramtext{تقسیم طبقه‌بندی‌شده با بذر ۴۲\\آموزش ۶۴ درصد، اعتبارسنجی ۱۶ درصد، آزمون ۲۰ درصد}};
\node[stage,below=of split] (clean) {\diagramtext{پاکسازی ثابت و کنترل هم‌پوشانی با اولویت آموزش\\۱۴۸٬۳۸۰ آموزش؛ ۳۷٬۰۸۶ اعتبارسنجی؛ ۴۶٬۳۴۲ آزمون}};
\node[stage,below=of clean] (cv) {\diagramtext{خط‌لولۀ واژه‌سازی، \en{Word2Vec} صدبعدی و طبقه‌بند\\اعتبارسنجی متقابل و تنظیم فراپارامتر فقط روی آموزش}};
\node[stage,below=of cv] (refit) {\diagramtext{برازش چهار مدل نهایی روی آموزش کامل\\واژگان و بازنمایی از اعتبارسنجی یا آزمون یاد گرفته نمی‌شوند}};
\node[stage,below=of refit] (select) {\diagramtext{انتخاب آستانۀ هر مدل و انتخاب مدل فقط بر اعتبارسنجی\\تثبیت انتخاب پیش از ارزیابی آزمون}};
\node[stage,below=of select] (test) {\diagramtext{ارزیابی آزمون برای مدل منتخب و گزارش دو آستانه\\ذخیرۀ معیارها، شمارش خطا، نمودارها و فرادادۀ اجرا}};
\foreach \a/\b in {raw/split,split/clean,clean/cv,cv/refit,refit/select,select/test}
  \draw[flow] (\a) -- (\b);
\end{tikzpicture}

\caption{جریان اصلاح‌شدۀ آزمایش متن؛ جداسازی ارزیابی پیش از یادگیری بازنمایی انجام می‌شود}
\label{fig:nlp-workflow-corrected}
\end{figure}
